\scanpage{22}
\[
\underset{\substack{\Vert\\\underline{\omega}_X}}{\det E}
\simeq \underset{\substack{\Vert\\M}}{\det F}\otimes L^{\otimes2},
\quad\text{donc }\deleted{\unsure{M\otimes}}\underline{\omega}_X\simeq M\otimes L^{\otimes2}
\]
donc (comme $M$ est de degré $0$) $2\deg L=\deg\underline{\omega}_X$, i.e. \deleted{deg}
\[
\deg L=g-1
\]
Considérons l'hom composé
\deleted{\diagram{\begin{tikzcd}\unsure{\underline{\omega}_X} \arrow[r] & \unsure{E}\end{tikzcd}}}
\diagram{\begin{tikzcd}
\underline{\omega} \arrow[r,hook] & E \arrow[r,"\simeq"{description}] & F\otimes L \arrow[r,two heads] & L
\end{tikzcd}}
d'où une section de $L\otimes\underline{\omega}^{-1}$, qui est de degré $-(g-1)<0$, donc \unsure{elle est} nulle, donc \deleted{cela se factorise par} $E\to L$ se factorise par $E/\underline{\omega}\simeq\mathcal O_X$, donc $L\simeq\mathcal O_X$ absurde car $\deg L\ne0$ !

Ainsi, le \deleted{faisceau} \inserted{groupe} structural du $\unsure{\xi|\zeta}$-fibré projectif associé à une conn. proj. ne peut se réduire à un groupe de Borel.

Montrons qu'il ne peut non plus se réduire à un sous-groupe \unsure{normalisateur} d'un tore maximal. En effet, si cela était le cas, alors par passage à un rev. étale, on se ramènerait au tore maximal, donc à un Borel, absurde (car le genre des revêtements \unsure{ser.} aussi $\geq2$ !). Le même argument montre :

\figure{\begin{tikzpicture}\draw (0,0)--(0,-2.6);\end{tikzpicture}}
\uline{\unsure{Pr.|Prop.}} Soit $X$ courbe projective lisse \inserted{conn.} de genre $g\geq2$ sur corps $k$ (de car nulle ?). Alors pour tte connexion projective sur $X$, on ne \unsure{peut} \deleted{$\unsure{\xi}$} réduire le groupe structural du $\unsure{\xi|\zeta}$-fibré principal de groupe $GP(1)_k=H$ associé, à un sous-groupe \unsure{algéb} $\unsure{H'|H}$ distinct de $H$....
\margin{!}
% archive foliation: 21
